% Standalone problem document, generated from the adjacent README.md.
% Compile with XeLaTeX. Mathematical content is maintained in Markdown.
\documentclass[11pt,a4paper]{article}
\usepackage[fontsize=10.5pt]{scrextend}
\usepackage[margin=22mm,headheight=15pt,headsep=9mm,footskip=11mm]{geometry}
\usepackage{fontspec}
\setmainfont{texgyrepagella}[Extension=.otf,UprightFont=*-regular,BoldFont=*-bold,ItalicFont=*-italic,BoldItalicFont=*-bolditalic]
\setsansfont{texgyreheros}[Extension=.otf,UprightFont=*-regular,BoldFont=*-bold,ItalicFont=*-italic,BoldItalicFont=*-bolditalic]
\setmonofont{texgyrecursor}[Extension=.otf,UprightFont=*-regular,BoldFont=*-bold,ItalicFont=*-italic,BoldItalicFont=*-bolditalic,Scale=MatchLowercase]
\usepackage{amsmath,amssymb}
\usepackage{unicode-math}
\setmathfont{texgyrepagella-math.otf}
\usepackage{xcolor,microtype,enumitem,needspace,fancyhdr}
\usepackage{bookmark}
\definecolor{ink}{HTML}{17344B}
\definecolor{accent}{HTML}{197D80}
\definecolor{muted}{HTML}{596773}
\hypersetup{colorlinks=true,linkcolor=accent,urlcolor=accent,pdftitle={MI-06 - Thompson-type
domination for the arithmetic symmetric
modulus},pdfauthor={Open Problems in Numerical Linear Algebra}}
\urlstyle{same}
\setlength{\parindent}{0pt}
\setlength{\parskip}{5pt plus 1pt minus 1pt}
\linespread{1.04}
\setlength{\emergencystretch}{3em}
\widowpenalty=10000
\clubpenalty=10000
\displaywidowpenalty=10000
\allowdisplaybreaks[1]
\setlist{nosep,leftmargin=1.5em}
\providecommand{\tightlist}{\setlength{\itemsep}{0pt}\setlength{\parskip}{0pt}}
\providecommand{\pandocbounded}[1]{#1}
\pagestyle{fancy}
\fancyhf{}
\fancyhead[L]{\sffamily\footnotesize\color{muted}OPEN PROBLEMS IN NUMERICAL LINEAR ALGEBRA}
\fancyhead[R]{\sffamily\bfseries\footnotesize\color{accent}MI-06}
\fancyfoot[L]{\parbox[b]{0.9\textwidth}{\sffamily\fontsize{8}{10}\selectfont\color{muted}Editorial ratings; a bounded search cannot certify that a problem remains unsolved.\\Literature check: 2026-09-12}}
\fancyfoot[R]{\sffamily\footnotesize\color{muted}\thepage}
\renewcommand{\headrulewidth}{0.3pt}
\renewcommand{\headrule}{\hbox to\headwidth{\color{accent}\leaders\hrule height \headrulewidth\hfill}}
\makeatletter
\renewcommand{\section}{\Needspace{5\baselineskip}\@startsection{section}{1}{0pt}{12pt plus 2pt minus 2pt}{4pt}{\sffamily\bfseries\large\color{ink}}}
\renewcommand{\subsection}{\Needspace{4\baselineskip}\@startsection{subsection}{2}{0pt}{9pt plus 1pt minus 1pt}{3pt}{\sffamily\bfseries\normalsize\color{ink}}}
\makeatother
\setcounter{secnumdepth}{0}
\begin{document}
{\sffamily\small\color{muted}Matrix inequalities and norms\par}
\vspace{3pt}
{\raggedright\hyphenpenalty=10000\exhyphenpenalty=10000\sffamily\bfseries\fontsize{21}{24}\selectfont\color{ink}Thompson-type
domination for the arithmetic symmetric modulus\par}
\vspace{7pt}
{\sffamily\small\textbf{Difficulty:} challenging\quad\textbf{Importance:} interesting
to specialist\par}
{\sffamily\small\bfseries\color{accent}Status: Lean verified\par}
\vspace{4pt}
{\color{accent}\hrule height 0.8pt}
\vspace{6pt}
\textbf{Rating rationale:} The arithmetic modulus lacks the useful order
structure of the quadratic modulus; a solution would advance a focused
family of operator triangle estimates.

\subsection{Lean proof and verification evidence ---
2026-09-12}\label{lean-proof-and-verification-evidence-2026-09-12}

\textbf{The original MI-06 assertion with factor \(\sqrt2\) is false,
with a complete Lean-verified counterexample.} At the rational pair

\[A=\begin{pmatrix}1&3/4&0\\0&0&0\\0&0&0\end{pmatrix},\qquad
B=\begin{pmatrix}-1&0&0\\0&0&0\\-3/4&0&0\end{pmatrix},\]

all six genuine matrix moduli are proved exactly. For every complex
unitary pair, a nonzero common orthogonal vector gives a left quadratic
form at least \((3/8)\|w\|_2^2\) and a right quadratic form at most
\((\sqrt2/4)\|w\|_2^2\). The strict scalar gap follows from the
kernel-checked LeanCert certificate \(2<9/4\). The
\href{https://github.com/sgstepaniants/OpenProblemsInNLA/tree/43b3dc65116633a68c32ec582fd093f3adc95597/matrix-inequalities-and-norms/MI-06/lean}{proof
at revision 43b3dc6} retains the full original complex-matrix and
all-unitary target.

\textbf{Mathematical counterexample and informal proof:} Matthew J.
Colbrook, Department of Applied Mathematics and Theoretical Physics,
University of Cambridge. \textbf{Lean formalization:} George
Stepaniants, Department of Computing and Mathematical Sciences,
California Institute of Technology, Pasadena, California, USA, with
AI-agent assistance.

The six checked declarations in
\href{https://github.com/sgstepaniants/OpenProblemsInNLA/blob/43b3dc65116633a68c32ec582fd093f3adc95597/matrix-inequalities-and-norms/MI-06/lean/Solution.lean}{Solution.lean}
are:

\begin{itemize}
\tightlist
\item
  \textit{NLA.MI06.modulus\_eq\_sqrt}: the genuine positive square-root
  modulus.
\item
  \textit{NLA.MI06.witness\_moduli}: six exact moduli, their averages
  and rank-one decompositions.
\item
  \textit{NLA.MI06.two\_vector\_orthogonal}: a nonzero common orthogonal
  vector in complex dimension three.
\item
  \textit{NLA.MI06.witness\_quadratic\_bounds}: homogeneous bounds for
  every complex unitary pair.
\item
  \textit{NLA.MI06.counterexample}: the fixed pair defeats all unitary
  choices.
\item
  \textit{NLA.MI06.not\_dominationConjecture}: negation of the complete
  original assertion.
\end{itemize}

Two independent agents reviewed the
\href{https://github.com/ajt60gaibb/OpenProblemsInNLA/blob/main/matrix-inequalities-and-norms/MI-06/lean/reviews}{frozen
statements and completed proof}.
\href{https://github.com/sgstepaniants/OpenProblemsInNLA/actions/runs/34711237623}{Linux
run 34711237623} matched all six exports with the sandboxed Comparator
and replayed them through Lean's default kernel. The
\href{https://github.com/ajt60gaibb/OpenProblemsInNLA/blob/main/matrix-inequalities-and-norms/MI-06/lean/verification/linux-2026-09-12}{original
artifacts and independent operational audit} retain all 117 source
hashes, isolation checks and rejection controls. All 52
\href{https://github.com/ajt60gaibb/OpenProblemsInNLA/blob/main/matrix-inequalities-and-norms/MI-06/lean/reviews/proof-build.log}{transitive
axiom reports} contain only \textit{propext}, \textit{Classical.choice}
and \textit{Quot.sound}. These are independent agent reviews, not
external human peer review.

The proof pins \textbf{Lean 4.33.1},
\href{https://github.com/alerad/leancert/tree/621a43d7cf21f87872392a01e874f2f1dbddc926}{LeanCert
621a43d} and
\href{https://github.com/leanprover-community/mathlib4/tree/0df444a360eaa60ab8c11dca51a86af692955474}{Mathlib
0df444a}. See the
\href{https://github.com/ajt60gaibb/OpenProblemsInNLA/blob/main/matrix-inequalities-and-norms/MI-06/lean/formalization.yaml}{manifest},
\href{https://github.com/ajt60gaibb/OpenProblemsInNLA/blob/main/matrix-inequalities-and-norms/MI-06/lean/lake-manifest.json}{dependency
pins} and
\href{https://github.com/ajt60gaibb/OpenProblemsInNLA/blob/main/matrix-inequalities-and-norms/MI-06/lean/NUMERICAL_TARGETS.md}{numerical
targets}. On a documented
\href{https://github.com/ajt60gaibb/OpenProblemsInNLA/blob/main/tools/lean/HARNESS.md}{non-root
Linux host}, reproduce from the verified revision with:

\begin{verbatim}
tools/lean/bootstrap.sh /absolute/path/to/nla-lean-tools
tools/lean/selftest.sh /absolute/path/to/nla-lean-tools
tools/lean/verify.sh \
  matrix-inequalities-and-norms/MI-06/lean \
  /absolute/path/to/nla-lean-tools
\end{verbatim}

The stronger no-finite-constant theorem in the original manuscript is
outside these six formal exports. The original informal proof and its
historical audit remain below.

\subsection{Resolution --- 2026-09-11}\label{resolution-2026-09-11}

\textbf{Negative result by Matthew J. Colbrook}, Department of Applied
Mathematics and Theoretical Physics, University of Cambridge.
\textbf{Independent proof review: PASS.}

No finite constant permits the proposed two-unitary Loewner-order
domination for the arithmetic symmetric modulus, already in dimension
three. A fixed rational example also refutes the proposed \(\sqrt2\)
constant. This concerns matrix order, not a separate norm triangle
inequality.

The exact target is resolved. The original statement and source evidence
are retained below; its former difficulty rating is historical.

\textbf{Primary manuscript:}
\href{https://github.com/ajt60gaibb/OpenProblemsInNLA/blob/main/matrix-inequalities-and-norms/MI-06/solution.pdf}{complete
proof PDF},
\href{https://github.com/ajt60gaibb/OpenProblemsInNLA/blob/main/matrix-inequalities-and-norms/MI-06/solution.tex}{standalone
TeX}, Theorem 1.1 and its proof;
\href{https://github.com/ajt60gaibb/OpenProblemsInNLA/blob/main/matrix-inequalities-and-norms/MI-06/solution.md}{authorship
and scope}. The
\href{https://github.com/ajt60gaibb/OpenProblemsInNLA/blob/main/references/colbrook-matrix-2026-09-11/verification/reviews/MI-06-review.md}{independent
review} checks the full original argument and records its hash. The
draft was AI-assisted; the 2026-09-11 review was independent agent
verification, without external human peer review or formal
certification. The later Lean verification above covers the complete
original factor-\(\sqrt2\) target.
\href{https://github.com/ajt60gaibb/OpenProblemsInNLA/blob/main/references/colbrook-matrix-2026-09-11/README.md}{Submission
record}.

\subsection{Problem statement}\label{problem-statement}

For \(X\in\mathbb C^{n\times n}\), define \(|X|=(X^*X)^{1/2}\) and

\[S(X)=\frac{|X|+|X^*|}{2}.\]

For every \(n\ge1\) and every \(X,Y\in\mathbb C^{n\times n}\), do there
exist unitary \(U,V\in\mathbb C^{n\times n}\) such that

\[S(X+Y)\preceq\sqrt2\bigl(U S(X)U^*+V S(Y)V^*\bigr)?\]

Here \(P\preceq Q\) means \(Q-P\) is positive semidefinite. The
conjectured universal factor \(\sqrt2\) is already known to be
necessary.

\subsection{Why it matters}\label{why-it-matters}

Symmetrizing the left and right polar factors removes their directional
preference. The question asks for a precise matrix-order bound for sums,
which would yield singular-value and norm consequences for these
positive representatives.

\newpage
\subsection{References}\label{references}

\begin{enumerate}
\def\labelenumi{\arabic{enumi}.}
\tightlist
\item
  T. Zhang, \emph{Operator symmetric moduli and sharp triangle
  inequalities}, arXiv:2603.01046v1 (1 March 2026), §1.2, Conjecture
  1.6. \href{https://arxiv.org/html/2603.01046}{Primary text}.
\item
  J.-C. Bourin and E.-Y. Lee, \emph{Some hybrid matrix triangle
  inequalities}, arXiv:2606.29188v1 (28 June 2026), introduction,
  Theorems 1.1 and 1.3. \href{https://arxiv.org/html/2606.29188}{Primary
  text}.
\end{enumerate}

\subsection{Status check --- 2026-09-10}\label{status-check-2026-09-10}

Both latest arXiv records remain v1. The June paper supplies related
estimates but no proof of the displayed two-unitary inequality. Searches
included \textit{Zhang Conjecture 1.6 symmetric moduli},
\textit{arithmetic symmetric modulus Thompson conjecture 2026}, and
the two titles with \textit{proof}. No resolution was located. The
analogous theorem for the quadratic symmetric modulus is proved and must
not be confused with this conjecture; a formula labeled Conjecture 3.13
in arXiv:2606.15624 appears to misidentify the modulus, so the draft
follows Zhang's original statement.

\textbf{Audit update (2026-09-10):} Rechecked Zhang Conjecture 1.6 and
Bourin--Lee Theorem 1.3, with searches for later symmetric-modulus
resolutions. Submajorization is weaker than the required two-unitary
positive-semidefinite domination. This is a bounded literature check,
not a proof that no solution exists.
\end{document}
